\section*{\S~3a Abwahl und Rücktritt}
\label{sec:abwahl_ruecktritt}
\addcontentsline{toc}{section}{\S~3a Abwahl und Rücktritt}

\begin{enumerate}[label=(\arabic*)]
    \item Die Abwahl wird im Folgenden geregelt.
    \begin{enumerate}[label=(\Alph*)]
      \item Eine Abwahl kann für jeden gewählten Posten der Unit durchgeführt werden und ist nur im Rahmen einer Mitgliederversammlung möglich.
      \item Eine Abwahl kann nur durch den S1-Feldwebel vorgeschlagen werden. Damit diese eingeleitet wird, benötig er zusätzlich die Zustimmung von 25 \% der Mitglieder. Ist der S1-Feldwebel vakant, so kann eine Abwahl nur durch den Gefreiten vom Dienst vorgeschlagen werden. Der Gefreite vom Dienst kann in diesem Fall auf Vorschlag von 25 \% der Mitglieder abgewählt werden. Die Abwahl des S1-Feldwebels kann in jedem Fall nur durch den Gefreiten vom Dienst vorgeschlagen werden.
      \item Die Abwahl ist erfolgreich, wenn sie eine Zweidrittelmehrheit erreicht.
      \item Bei erfolgreicher Abwahl scheidet der Amtsinhaber sofort aus seinem Amt aus und eine Neuwahl muss innerhalb von 14 Tagen durchgeführt werden. Der Posten bleibt so lange vakant
      \item Falls eine Abwahl scheitert, kann frühestens nach drei Monaten erneut ein Antrag auf Abwahl des Amtsinhabers gestellt werden.
    \end{enumerate}
    \item Der Rücktritt wird im Folgenden geregelt.
      \begin{enumerate}[label=(\Alph*)]
      	\item Ein Amtsinhaber kann zu jedem Zeitpunkt freiwillig zurücktreten.
	\item Die Neuwahl muss innerhalb von 14 Tagen stattfinden.
    \end{enumerate}
\end{enumerate}